\subsection*{初始化空槽事件}

循环开始前，需要先触发空槽对应的 \co@empty@ 事件：

\begin{lstlisting}[language=croqtile,caption={引导协议}]
foreach {s} in [STAGES] {
  trigger empty[s];
}
\end{lstlisting}

事件初始为未触发状态。如果没有这一步，生产者第一次等待空槽时就无法继续。初始化表示这些槽位已经可以使用。

完整的双缓冲和多阶段流水线示例见 3.8.5 节。先理解事件的触发与等待关系，再阅读这些组合用法。
